\section{Thesis outline}
The goal of this thesis is to design and implement a NoC router for efficient data movement in a small mesh-based communication system. The router is intended to support packet-based communication between processing elements and to provide basic routing, virtual-channel flow control, and pipelined forwarding. A 3×3 mesh network is used as the target structure, and FPGA implementation is considered as the hardware platform for evaluating the design.

The thesis work is organized as follows:

\begin{itemize}
    \item \textbf{Chapter 2: Background.} 
    Introduces the key concepts used in this work, including chiplet-based systems, 
    Network-on-Chip, mesh topology, packet switching, routing, virtual channels, 
    and flow control.

    \item \textbf{Chapter 3: Design and Implementation.} 
    Describes the proposed router architecture and its main design components. 
    This chapter explains the flit format, routing algorithm, allocation methods, 
    pipeline organization, and 3$\times$3 mesh integration.

    \item \textbf{Chapter 4: Verification and Results.} 
    Presents the verification plan and FPGA-oriented implementation method. 
    It describes single-router verification, mesh-level verification, test cases, 
    and the use of UART for test control and result reporting.

    \item \textbf{Chapter 5: Conclusion and Future Work.} 
    Summarizes the thesis work and discusses possible future improvements, such as 
    extending the design to different network topologies, exploring adaptive routing 
    algorithms, improving arbitration schemes, and integrating the NoC with real 
    processing elements or accelerator modules.
\end{itemize}